%%%% Time-stamp: <2013-02-25 10:31:01 vk>


\chapter*{Kurzfassung /Abstract }
\label{cha:abstract}

Eine Kurzfassung ist in deutscher sowie ein Abstract in englischer Sprache mit je maximal einer A4-Seite zu erstellen. Die Beschreibung sollte wesentliche Aspekte des Projektes in technischer Hinsicht beschreiben. Die Zielgruppe der Kurzbeschreibung sind auch Nicht-Techniker! Viele Leser lesen oft nur diese Seite. \\ \\

In einer Welt in der digitale Technik und Systeme immer mehr an Bedeutung gewinnen, haben wir uns bei dieser Diplomarbeit damit beschäftigt, wie wir diese Technik uns zu Nutze machen können um Menschen den Einstieg ins Schachspiel zu vereinfachen. Ziel dieser Diplomarbeit  war die Entwicklung und Umsetzung eines Einstiegsfreundlichen Schachbretts, welches den Spieler bei der Richtigen Zugwahl unterstützt, visuell zeigt welche Züge gezogen werden dürfen und mittels Spielynalyse, den Spieler auf Fehler hinweisen kann. Ein zentrale Bestandteile der Diplomarbeit war, die Entwicklung der Hardware und das Design des Schachbretts, Ein weiterer wichtiger Aufgabenbereich der Arbeit war, die Programmierung der entwickelten Hardeware und die softwareseitige Einbindung der Bedienelemente. Der letzte große Aufgabenbereich bestand darin die lokal gespeicherten Daten mit einem Server zu kommuniezieren. Auf diesem Server kann der Nutzer sein Profil anlegen und seine zuvor gespielten Partien analysieren.    \\ \\


This thesis deals with the various questions concerning learning for adults – with the aim to describe learning cultures which support the concept of live-long learning (LLL). The learning ability of adults and the various motives which lead to adults learning are the starting point of this thesis. The following analysis on self-directed learning as well as the resulting new attribution of roles and tasks which arise for learners, trainers and institutions in adult education, shall demonstrate first possibilities to contribute to the implementation of the concept of LLL. In addition, the role of information and communication technologies in the framework of LLL will be closer described in context of self-directed learning processes of adults as the opening of new forms of communication and co-operation independent of location and time between learners as well as between learners and tutors gains more importance. Finally the topic of visualisation, validation and recognition of informal and non-formal learning and their contribution to LLL is discussed. \\

Gliederung des Abstract in \textbf{Thema}, \textbf{Ausgangspunk}, \textbf{Kurzbeschreibung}, \textbf{Zielsetzung}.  

\subparagraph{Projektergebnis}

Allgemeine Beschreibung, was vom Projektziel umgesetzt wurde, in einigen kurzen Sätzen. Optional Hinweise auf Erweiterungen. Gut machen sich in diesem Kapitel auch Bilder vom Gerät (HW) bzw. Screenshots (SW).
Liste aller im Pflichtenheft aufgeführten Anforderungen, die nur teilweise oder gar nicht umgesetzt wurden (mit Begründungen).